\documentclass[11pt,a4paper]{article}
\usepackage[margin=25mm]{geometry}
\usepackage{amsmath,amssymb,amsthm}
\usepackage{hyperref}

\newtheorem{theorem}{Theorem}[section]
\newtheorem{proposition}[theorem]{Proposition}
\newtheorem{lemma}[theorem]{Lemma}
% Choose a number before each manually numbered theorem.
\newcommand{\manualnumber}{}
\newtheorem*{manualtheorem}{Theorem \manualnumber}
\theoremstyle{definition}
\newtheorem{definition}[theorem]{Definition}
\newtheorem{example}[theorem]{Example}

\title{Algebra 1}
\author{roger}
\date{}

\begin{document}
\maketitle
\tableofcontents

\section{Lecture 1: Topic}\label{sec:lecture-1}
\noindent\textbf{Date:} \hfill \textbf{Key idea:}

\subsection{Definitions and results}
% Record notation, definitions, and theorem statements.

\subsection{Examples and proofs}
% Include a worked example and any important proof steps.

\subsection{Questions to revisit}
% Note any gaps to resolve after the lecture.

\subsection{Revision summary}
% Add a few sentences after reviewing the lecture.

\subsection{Steinitz exchange lemma}
\renewcommand{\manualnumber}{2.15}
\begin{manualtheorem}
If a subspace $S$ has a spanning sequence of length $m$, then any
sequence of more than $m$ elements of $S$ is linearly dependent.
\end{manualtheorem}

\begin{proof}[Proof (incomplete)]
% Omit the end-of-proof square until the argument is finished.
\renewcommand{\qedsymbol}{}
Let $s_1,\ldots,s_m$ be a spanning sequence for $S$.
Let $n>m$, and let $v_1,\ldots,v_n$ be any sequence of elements of $S$.
We must show that $v_1,\ldots,v_n$ are linearly dependent.

\emph{To be completed.}
% TODO: add the remaining proof steps, then remove "(incomplete)",
% the line changing \qedsymbol, and "To be completed."
\end{proof}

\renewcommand{\manualnumber}{2.17}
\begin{manualtheorem}
Any two bases of a subspace have the same size.
\end{manualtheorem}


\noindent This theorem can be proven easily
\\
about 


\section{Subspaces}
Applications of linear algebra to geometry work best when the shapes involved 
are straight lines, plaes and their generalization s in higher dimentions.

\begin{definition}
    (a) A set U of elements of $F^n$ is called closed under addition if for 
    every x,y in U their sum x+y belongs to U \\ \\
    (b) A set U of elements of $F^n$ is called closed under scaler multiplication if 
    for every x in U and every number a in F, the scalar multiple ax belongs to U.\\\\
    (c) A subset U of $F^n$ is called a subspace of $F^n$ if it is non-empty, closed under addition,
    and closed under scalar multiplication. We write $U<=F^n$ to mean that U is a subspace of $F^N$
\end{definition}

Sometimes one subspace U is contained in another subspace V, that is every element of U also belongs to V.
We write as $U<=V$ and U is a \textbf{proper} subspace if $U<=V$ and $U!=V$

Another way is 

\begin{definition}
  Let A be an $m*n$ matrix with entries in $F$.
  
  (a) The \textbf{kernal} of A, wrriten $N(A)$, is ${v\in F^n:Av=0_m}$

  (b) The \textbf{column space} of A, written C(A),is ${AV:V\in F^n}$
\end{definition}

\begin{lemma}[Kernals and column spaces are subspaces]
    Let A be an $m*n$ matrix with entries in $F$. Then $N(A)\leq \mathbb{F}^n$ and $C(A)<=\mathbb{F}^m$.

\end{lemma}

\begin{proof}
  We chekc 3 subspace conditions(non-empty,closed under addition and closure under scalar multiplication).

  (a) Kernal. $A0_n=0_m$, so $0_n \in N(A)$ which is therefore non-empty. Let u,v $\in N(A)$ and let $a\in \mathbb{F}$. Then
\begin{align*}
A(u+v) &= Au+Av\\
&=0_m+0_n\\
&=0_m 
\end{align*}


\end{proof}

section



















\end{document}
